\ea\label{ex:text28-1}
O haopi a hena na iana vai to teitei batana paano kasuana.\\
\gll o	haopi	a	hena	n=a	iana	vai	to	tei{$\sim$}tei	bata=na	paa-n=o	kasuana\\
     \textsc{art}3.\textsc{sg}	wrasse	\textsc{art}2.\textsc{sg}	name	3\textsc{sg}.\textsc{poss}=\textsc{art}2.\textsc{sg}	fish	\textsc{dem}6	\textsc{rel}	\textsc{red}{$\sim$}stay	\textsc{cont}=3\textsc{sg}.\textsc{ipfv}	under-3\textsc{sg}.\textsc{poss}=\textsc{art}3.\textsc{sg}	ground\\
\glt `\textit{Haopi} (Wrasse) is the name of a fish that stays underground.'
\z

\ea\label{ex:text28-2}
Na tei me nana bono buaku o kaku haopi.\\
\gll na	tei	me	nana	bono	buaku	o	kaku	haopi\\
     \textsc{tm}2	stay	with	3\textsc{sg}.\textsc{ipfv}	\textsc{obj}.\textsc{art}3.\textsc{sg}	two	\textsc{art}3.\textsc{sg}	kind	wrasse\\
\glt `There are two kinds of wrasse.'\footnote{lit. `(It) is there with two kinds of wrasse.'}
\z

\ea\label{ex:text28-3}
O vuaba vai o beera na dao raara bono mataberam.\\
\gll o	vuaba	vai	o	beera	na	dao	ra=ara	bono	mataberam\\
     \textsc{art}3.\textsc{sg}	one	\textsc{dem}6	\textsc{art}3.\textsc{sg}	big	\textsc{tm}2	call	1\textsc{pl}.\textsc{in}.\textsc{ipfv}=1\textsc{pl}.\textsc{in}.\textsc{pron}	\textsc{obj}.\textsc{art}3.\textsc{sg}	Rockmover.Wrasse\\
\glt `The big one is called \textit{mataberam} (Rockmover Wrasse).'
\z

\ea\label{ex:text28-4}
O vuaba vai o varutaa na dao raara bono haopi suuna.\\
\gll o	vuaba	vai	o	va-rutaa	na	dao	ra=ara	bono	haopi	suuna\\
     \textsc{art}3.\textsc{sg}	one	\textsc{dem}6	\textsc{art}3.\textsc{sg}	\textsc{caus}-small	\textsc{tm}2	call	1\textsc{pl}.\textsc{in}.\textsc{ipfv}=1\textsc{pl}.\textsc{in}.\textsc{pron}	\textsc{obj}.\textsc{art}3.\textsc{sg}	wrasse	genuine\\
\glt `The smaller one is called \textit{haopi} \textit{suuna} (genuine wrasse).'
\z

\ea\label{ex:text28-5}
Evehee o hena beera o haopi.\\
\gll evehee	o	hena	beera	o	haopi\\
     but	\textsc{art}3.\textsc{sg}	name	general	\textsc{art}3.\textsc{sg}	wrasse\\
\glt `But the general name is wrasse.'
\z

\ea\label{ex:text28-6}
Beara tau vagana haopi, eara repaa tara rakerake bata o sii vahara vasu beori ravun vaa tea vaakororovi rori.\\
\gll be=ara	tau	vagana	haopi	e-ara	re=paa	tara	rake{$\sim$}rake	bata	o	sii	vahara	vasu	be=ori	ravun	vaa	te=a	vaakororovi	rori\\
     when=1\textsc{pl}.\textsc{in}.\textsc{pron}	\textsc{tm}5	fish	wrasse	\textsc{h}-1\textsc{pl}.\textsc{in}.\textsc{pron}	\textsc{consec}=\textsc{tm}3	see	\textsc{red}{$\sim$}search	\textsc{cont}	\textsc{art}2.\textsc{pl}	\textsc{dim}	small.\textsc{pl}	stone	when=3\textsc{pl}.\textsc{pron}	be.heaped.up	\textsc{lnk}	\textsc{prep}=\textsc{art}2.\textsc{sg}	circle	3\textsc{pl}.\textsc{ipfv}\\
\glt `When we are about to catch wrasses, we search for little stones being heaped up in a circle.'
\z

\ea\label{ex:text28-7}
Beara tara ahe eve, eara toro huu hopo ni nao a sinaka,\\
\gll be=ara	tara	ahe	eve	e-ara	toro	huu	hopo	ni	nao	a	sinaka\\
     when=1\textsc{pl}.\textsc{in}.\textsc{pron}	see	recognize	3\textsc{sg}.\textsc{pron}	\textsc{h}-1\textsc{pl}.\textsc{in}.\textsc{pron}	\textsc{tm}7	push	enter	\textsc{appl}	thither	\textsc{art}2.\textsc{sg}	fishing.stick\\
\glt `When we spot it, we must push the fishing-stick into it,'
\z

\ea\label{ex:text28-8}
ei pasi pahin tahi ahe koa vuara.\\
\gll ei	pasi	pahin	tahi	ahe	koa	vu-ara\\
     \textsc{dem}2	\textsc{tm}4	immediately	spear	recognize	just	\textsc{pfv}-1\textsc{pl}.\textsc{in}.\textsc{pron}\\
\glt `this (means) we will immediately spear it.'
\z

\ea\label{ex:text28-9}
Eara repaa kasi, are kasi antee beara kasi ahe komana e.\\
\gll e-ara	re=paa	kasi	a=re	kasi	antee	be=ara	kasi	ahe	komana	e\\
     \textsc{h}-1\textsc{pl}.\textsc{in}.\textsc{pron}	\textsc{consec}=\textsc{tm}3	dig	1\textsc{pl}.\textsc{in}.\textsc{pron}=\textsc{consec}	dig	stop	when=1\textsc{pl}.\textsc{in}.\textsc{pron}	dig	recognize	self	3\textsc{sg}.\textsc{pron}\\
\glt `Then we will dig and dig until we recognize it.'\footnote{lit. `Then we dig, then we dig (and) stop, when dig (and) recognize ...'}
\z

\ea\label{ex:text28-10}
Ato nie teo buo ketaa nae, eara repaa rahi tavusu e.\\
\gll ato	ni=e	te=o	buo	ketaa	nae	e-ara	re=paa	rahi	tavusu	e\\
     hold	\textsc{appl}=3\textsc{sg}.\textsc{pron}	\textsc{prep}=\textsc{art}3.\textsc{sg}	two3.\textsc{sg}	gill	3\textsc{sg}.\textsc{poss}.\textsc{pron}	\textsc{h}-1\textsc{pl}.\textsc{in}.\textsc{pron}	\textsc{consec}=\textsc{tm}3	pull	come.out	3\textsc{sg}.\textsc{pron}\\
\glt `(We) hold it on its two gills, and then we pull it out.'
\z

\ea\label{ex:text28-11}
Ei asun vamate vai ara.\\
\gll ei	asun	va-mate	vai	ara\\
     \textsc{dem}2	hit	\textsc{advr}-dead	now	1\textsc{pl}.\textsc{in}.\textsc{pron}\\
\glt `Now we kill it.'
\z

\ea\label{ex:text28-12}
Evehee na teitei batana a sii vapara nae.\\
\gll evehee	na	tei{$\sim$}tei	bata=na	a	sii	vapara	nae\\
     but	\textsc{tm}2	\textsc{red}{$\sim$}be.there	\textsc{cont}=3\textsc{sg}.\textsc{ipfv}	\textsc{art}2.\textsc{sg}	\textsc{dim}	partner	3\textsc{sg}.\textsc{poss}.\textsc{pron}\\
\glt `But usually there is also its partner.'
\z

\ea\label{ex:text28-13}
Goe tea nao vananaoana vira bean asun vakavara o peho.\\
\gll goe	tea	nao	va-nanaoana	vira	be=an	asun	va-kavara	o	peho\\
     don't	\textsc{compl}1	go	\textsc{advr}-far.away	too	when=2\textsc{sg}.\textsc{pron}	hit	\textsc{advr}-complete	\textsc{art}3.\textsc{sg}	one\\
\glt `Don't go too far away when you have killed one.'
\z

\ea\label{ex:text28-14}
Tara rakerake vaha koa o meho rihinae.\\
\gll tara	rake{$\sim$}rake	vaha	koa	o	meho	rihi-nae\\
     see	\textsc{red}{$\sim$}search	again	just	\textsc{art}3.\textsc{sg}	other3.\textsc{sg}	proximity-3\textsc{sg}.\textsc{poss}.\textsc{pron}\\
\glt `Just search for the other one nearby.'\footnote{lit. `... the other one in its proximity'}
\z

\ea\label{ex:text28-15}
Ean paa kasi e.\\
\gll e-an	paa	kasi	e\\
     \textsc{h}-2\textsc{sg}.\textsc{pron}	\textsc{tm}3	dig	3\textsc{sg}.\textsc{pron}\\
\glt `You dig it out.'
\z

\ea\label{ex:text28-16}
Eara samin doana haara, na pahin kasi ona bata koa raara.\\
\gll e-ara	sa-min	doana	haa=ra	na	pahin	kasi	ona	bata	koa	ra=ara\\
     \textsc{h}-1\textsc{pl}.\textsc{in}.\textsc{pron}	\textsc{neg}1-ever	fail	\textsc{neg}2=1\textsc{pl}.\textsc{in}.\textsc{ipfv}	\textsc{tm}2	immediately	dig	not.miss	\textsc{cont}	just	1\textsc{pl}.\textsc{in}.\textsc{ipfv}=1\textsc{pl}.\textsc{in}.\textsc{pron}\\
\glt `We never fail, we immediately dig (and) catch (it).'
\z

\ea\label{ex:text28-17}
Eara toro matatopo bata koara to mene rosin tavusu maa.\\
\gll e-ara	toro	matatopo	bata	koa=ra	to	mene	rosin	tavusu	maa\\
     \textsc{h}-1\textsc{pl}.\textsc{in}.\textsc{pron}	\textsc{tm}7	be.prepared.for	\textsc{cont}	just=1\textsc{pl}.\textsc{in}.\textsc{ipfv}	\textsc{rel}	\textsc{tm}6	escape	come.out	hither\\
\glt `We must be prepared that it might come out and escape.'
\z

\ea\label{ex:text28-18}
Kikira vakikis bata koa ni nom eve.\\
\gll kiikira	va-kikis	bata	koa	ni	nom	eve\\
     watch	\textsc{advr}-strong	\textsc{cont}	just	\textsc{appl}	\textsc{ipfv}	3\textsc{sg}.\textsc{pron}\\
\glt `Just keep carefully watching it .'
\z

\ea\label{ex:text28-19}
Eara mee toro goe tea ato o haopi en vihi nae ei kou to mene dohi ore paa rosin.\\
\gll e-ara	mee	toro	goe	tea	ato	o	haopi	en	vihi	nae	ei~kou	to	mene	dohi	o=re=paa	rosin\\
     \textsc{h}-1\textsc{pl}.\textsc{in}.\textsc{pron}	also	\textsc{tm}7	do.not	\textsc{compl}1	hold	\textsc{art}3.\textsc{sg}	wrasse	\textsc{dem}3	tail	3\textsc{sg}.\textsc{poss}.\textsc{pron}	because	\textsc{rel}	\textsc{tm}6	slip	3.\textsc{pron}=\textsc{consec}=\textsc{tm}3	escape\\
\glt `Also, we must not hold the wrasse on its tail because it might slip (out of your hands) and escape.'\footnote{lit. `... because (it might be one) that otherwise slips and then escapes'}
\z

\ea\label{ex:text28-20}
Evei o sii vahutate vaa tea vagana haopi.\\
\gll eve=i	o	sii	vahutate	vaa	tea	vagana	haopi\\
     3\textsc{sg}.\textsc{pron}=\textsc{dem}1	\textsc{art}3.\textsc{sg}	\textsc{dim}	description	\textsc{lnk}	\textsc{compl}1	fish	wrasse\\
\glt `This is the little description of fishing wrasses.'
\z
